% ===== Джерела до розділу 14. Ключі з префіксом ch14: =====
% Позиції ядра (bib-core.tex) тут НЕ дублюються.

% --- Промислові системи керування (ICS/SCADA) ---
\bibitem{ch14:stouffer} Stouffer K., Pillitteri V., Lightman S., Abrams M., Hahn A. Guide to Industrial Control Systems (ICS) Security. NIST Special Publication 800-82, Revision~2. Gaithersburg: NIST, 2015. DOI: 10.6028/NIST.SP.800-82r2.

\bibitem{ch14:langner} Langner R. Stuxnet: Dissecting a Cyberwarfare Weapon. \emph{IEEE Security \& Privacy}. 2011. Vol.~9, №~3. P.~49--51. DOI: 10.1109/MSP.2011.67.

\bibitem{ch14:industroyer} Cherepanov A. Win32/Industroyer: A New Threat for Industrial Control Systems. ESET White Paper, 2017.

\bibitem{ch14:iec62443} International Electrotechnical Commission. IEC~62443: Industrial Communication Networks --- Network and System Security. Geneva: IEC, 2018--2021.

% --- Інтернет речей (IoT) ---
\bibitem{ch14:antonakakis} Antonakakis M., April T., Bailey M., et al. Understanding the Mirai Botnet. \emph{Proceedings of the 26th USENIX Security Symposium}. USENIX Association, 2017. P.~1093--1110.

\bibitem{ch14:iot-owasp} OWASP Foundation. OWASP Internet of Things (IoT) Top 10. OWASP Project, 2018.

% --- Мобільні платформи ---
\bibitem{ch14:enck} Enck W., Octeau D., McDaniel P., Chaudhuri S. A Study of Android Application Security. \emph{Proceedings of the 20th USENIX Security Symposium}. USENIX Association, 2011. P.~21--36.

\bibitem{ch14:android} Google. Android Security Model and Application Sandbox --- Android Open Source Project Documentation. 2024. URL: \url{https://source.android.com/docs/security}.

% --- Хмарні середовища ---
\bibitem{ch14:cloud-nist} Mell P., Grance T. The NIST Definition of Cloud Computing. NIST Special Publication 800-145. Gaithersburg: NIST, 2011.

\bibitem{ch14:container} Souppaya M., Morello J., Scarfone K. Application Container Security Guide. NIST Special Publication 800-190. Gaithersburg: NIST, 2017.

% --- Українські джерела ---
\bibitem{ch14:ics-ua} Захист автоматизованих систем керування технологічними процесами об'єктів критичної інфраструктури України від кіберзагроз. \emph{Кібербезпека: освіта, наука, техніка}. 2024. URL: \url{https://csecurity.kubg.edu.ua/index.php/journal/article/view/788}.

\bibitem{ch14:iot-ua} Моделі та методи виявлення шкідливої активності у середовищах Інтернету речей. \emph{Захист інформації}. 2023. Т.~25, №~4.
